\documentclass[11pt]{article}
\usepackage[utf8]{inputenc}
\usepackage{amsmath,amssymb,amsthm}
\usepackage[margin=1in]{geometry}
\usepackage{hyperref}
\usepackage{xcolor}

\newcommand{\tideref}[2][]{\href{https://therisensea.org/tides/#2/}{\texttt{#2}}}
\emergencystretch=1.5em

\title{Degree-$k$ radial Taylor coefficients\\
(tensor programme J4)}
\author{Tide loop (Claude Fable 5, with GPT-5.6 Sol deliberation)}
\date{2026-08-09}

\begin{document}
\maketitle

\begin{abstract}
\noindent The analytic input for the tensor programme's integration
stage: for two $C^k$ losses whose derivative tensors at the origin
agree below order $k$,
\[
\frac{L_1(qx) - L_2(qx)}{q^k}
\;\xrightarrow[q \to 0^+]{}\;
\Delta_k(x)
:= \frac{1}{k!}\Bigl(D^kL_1(0) - D^kL_2(0)\Bigr)[x,\dots,x],
\]
with $\Delta_k$ homogeneous of degree $k$ --- exactly the certificate
the J2 rigidity engine consumes. The technical yield is a one-lemma
replacement for the quadratic stage's derivative plumbing: composing
\texttt{iteratedFDeriv} with the ray's continuous linear map gives
the ray derivatives at \emph{all} orders at once, where stage H3a
had assembled orders one and two by hand with the chain rule.
\end{abstract}

\section{Setting}

Stages J2 (covariance rigidity) and J3 (polarization) supply the
recovery logic; what they consume are diagonal homogeneous functions
arising as limits of rescaled loss differences. The scoping consult's
J4 prescribes the pairwise form directly --- comparing two losses
with matched lower jets --- rather than a standalone Taylor expansion,
because the recovery induction only ever sees differences.

\section{How the proof works}

The ray restriction $q \mapsto L(qx)$ is $L$ composed with the
continuous linear map $t \mapsto tx$, so Mathlib's
composition rule for iterated derivatives\footnote{%
\texttt{ContinuousLinearMap.iteratedFDeriv\_comp\_right} with
\texttt{iteratedDeriv\_eq\_iteratedFDeriv}.} evaluates every ray
derivative at the origin as the corresponding Fr\'echet diagonal in
one step. For the difference of two losses with matched jets below
$k$, every Taylor coefficient of the ray below order $k$ vanishes
(\texttt{iteratedFDeriv\_sub\_apply} plus the hypothesis), the
order-$k$ coefficient is $k!\,\Delta_k(x)$, and the one-dimensional
Peano remainder plus the established division-and-congruence
assembly gives the limit. Homogeneity of $\Delta_k$ is
\texttt{MultilinearMap.map\_smul\_univ}. A pleasant surprise: the
positivity hypothesis $0 < k$ turned out unnecessary --- at $k = 0$
the statement degenerates to continuity and the same proof closes it.

\section{Where it stalled and how it moved}

Nowhere: a two-repair diagnostic pass (one dead trailing tactic, one
unused hypothesis), the arc's cleanest tide after H6. The numerical
check (cubic difference at $d = 2$, linear-in-$q$ convergence to
$\Delta_3(x)$) ran before the Lean was written.

\section{Mathlib lookup versus new mathematics}

Lookup: the iterated-derivative composition rule, the difference
rule, and the multilinear scaling identity. New (project-level):
the packaged pairwise limit and the observation that the ray-CLM
composition subsumes the quadratic stage's hand-rolled derivative
computations.

\section{What this opens}

J5, the programme's flagged largest stage: integrating this tide's
pointwise limit against the posterior at the rate-sensitive scale
$q^{k-2}$, which the consult warns needs a local/tail split beyond
the H4 dominated-convergence pattern (a fixed Gaussian dominator no
longer suffices after dividing by $q^{k-2}$). Its hypothesis
structure (`HigherLaplaceDomain`) will extend the H4 package with
bounded $k$-th derivatives near the origin and a tail gap.

\end{document}
